\chapter{Simulating ARA's Performance with \texttt{PyREx} \label{app:veff_pyrex}}

  We simulate initial neutrinos at discrete energies spaced logarithmically from $3\times10^{17}$~eV to $10^{21}$~eV entering Earth from all angles and in just one array-wide configuration with all stations active.
  We simulate all subsequent interactions and cascades induced by outgoing leptons with energy greater than $10^{16}$ eV.
  All resulting particle cascades that occur in a cylinder with a 15~km radius are passed into a signal simulation package where ARA's triggering is simulated.

  The package that simulates the propagation of initial neutrinos, their interactions, and the propagations and interactions of all subsequent particles is NuLeptonSim \citep{nuleptonsim}.
  The model of Earth simulated follows the PREM model with a 2.8 kilometer layer of ice replacing the outermost 2.8 kilometers of Earth \citep{prem}.
  Depending on their energy and flavor, between 1,000 and 10,000 primary neutrinos were injected along 100,000 unique trajectories that are evenly spread across Earth and that intersect with ARA's viewing region.
  The distance traversed by each neutrino along its trajectory is randomly chosen according to the neutrino's cross section, which determines its interaction vertex.
  All subsequent outgoing neutrinos are propagated the same way.
  NuLeptonSim is similar to the existing PROPOSAL framework but has a faster implementation.

  The outgoing muons and taus are propagated differently: particle cascades along muon and tau tracks are determined stochastically (in position and energy) along the charged lepton's trajectory until the particle randomly decays or has less than $10^{16}$ eV of energy remaining.
  Outgoing neutrinos and charged leptons from muon and tau decays are also tracked, allowing us to incorporate $\nu_{\tau}\to\tau\to\nu_{\tau}$ regeneration in our study.
  By simulating neutrinos all over Earth, we are also able to include particles that are created in interactions and decay far away but eventually travel into our trigger volume (traditionally, initial neutrino cascades are only simulated within the trigger volume).
  All particle interactions and energy depositions in the 15-kilometer-radius trigger volume are characterized by location, energy, direction of travel, inelasticity, and interaction channel and passed to the next software.

  Python for Radio Experiments, \texttt{PyREx}, simulates the radio signal emitted by the nearby cascades, then simulates the waveforms generated by ARA antennas upon receipt of those signals by all five of our stations\footnote{https://github.com/abigailbishop/pyrex}.
  \texttt{PyREx} does this by first ray tracing and attenuating the Askaryan signals from the position of an event to each radio antenna that can observe the cascade.
  The electric field of the signal at each antenna is then convolved with the antenna's response function and modified through electronics filters until a voltage trace is achieved.
  ARA's trigger is formed when the signal from multiple antennas exceeds a threshold determined by the power generated by a signal in comparison to the power from background noise.
  So, the pure signal voltage waveform and a pure noise voltage waveform are both convolved through a tunnel diode to get power traces for each.
  For one station, if three or more antennas of the same polarization observe that the power generated from a cascade's signal is more than six times greater than the power generated by noise, the station triggers.
  This is the classic ARA trigger; we do not consider Phased Array triggers in this study as this functionality is not built into \texttt{PyREx} yet.

  \subsection{Muons and Taus}

    Given ARA's multiple stations spread across multiple cubic kilometers of South Pole Ice, cascades from outgoing leptons are expected to have a considerable effect on our sensitivity.
    At ultra-high energies, taus have a decay length of tens-of-meters allowing us to observe tau tracks as well as tau decays \citep{nuradiomc2}.
    This allows us to study a few interesting event topologies such as multi-particle events, muon and tau tracks, and muon and tau decays explained below and illustrated in Figure~\ref{fig:eventtopos}.
    
    \begin{enumerate} 
    

      \item Multi-cascade events are those where showers from one or more particles trigger our array.
      For example, a primary neutrino creates an initial cascade and the outgoing lepton generates secondary particle cascades in the ice that trigger the same ARA station.
      Another would be multiple stochastic losses from the same muon triggering 2 ARA stations.
      (Figure~\ref{fig:eventtopos}a)

      \item For tracks, a muon or tau resulting from a charged current neutrino interaction 
      %, a charged lepton decay, or a $W^-$ decay
      travels through our array's triggering region.
      As the muon or tau travels, it stochastically sheds energy, initiating particle cascades that emit Askaryan Radiation which may be detected in more than one station.
      Each deposit may be tens or hundreds of meters apart and each deposit, if observed, should have similar angular reconstruction and occur in time coincidence.
      This allows the observed energy depositions to be associated with each other.
      (Figure~\ref{fig:eventtopos}b)
      

      \item The last interesting event topology is one in which a tau or muon decays and the associated particle cascades are observed by our array.
      (Figure~\ref{fig:eventtopos}c)
      
    \end{enumerate}

    \begin{figure}
      \centering
      \includegraphics[
        width=0.6\linewidth,
        alt={x}
      ]{sections/araveff/images/icrc23/plots-1s3a6p-general-Aeffs-prim}
      \caption{
        Effective Area of all events involving an initial neutrino cascade (in black) compared to previous ARA initial cascade results (in red) \citep{ara_a23_10mo}.
        The blue lines correspond to the effective area for different flavors of neutrinos.
      }
      \label{fig:EAbyflavor}
    \end{figure}

    \begin{figure}
      \centering
      \includegraphics[
        width=0.6\linewidth,
        alt={x}
      ]{sections/araveff/images/icrc23/plots-1s3a6p-general-Aeffs-byNTriggered-all_flavor-LOFalse}
      \caption{
        Percent of effective area composed by events where $only$ an initial cascade trigger (in blue) and events where two or more cascades from the same event trigger (gray lines) one or more ARA stations.
      }
      \label{fig:EAbyNTriggers}
    \end{figure}
    
    \begin{figure}
      \centering
      \includegraphics[
        width=0.6\linewidth,
        alt={x}
      ]{sections/araveff/images/icrc23/plots-1s3a6p-general-Aeffs-allParticle}
      \caption{
        Percent of the Effective Area made up of: 
          all events including an initial cascade trigger (in blue, including events where cascades from an outgoing lepton triggered as well); 
          events where only cascades from taus, muons, and outgoing neutrinos trigger (in yellow);
          events where only outgoing neutrinos triggered (in green, all data points are between 0.5\% and 0.9\%).
      }
      \label{fig:EAbyparticles}
    \end{figure}
    
    \begin{figure}
      \centering
      \includegraphics[
        width=0.6\linewidth,
        alt={x}
      ]{sections/araveff/images/icrc23/plots-1s3a6p-coincident-Aeffs-percentAeffTot-all_flavor}
      \caption{
        Percent of total ARA effective area from: 
          all multi-station events (in black), 
          initial neutrino cascades (in blue), 
          a single muon or tau cascade (from a track or decay) triggering multiple stations (in solid yellow), 
          multiple muon and tau cascades triggering different stations (in dash-dot yellow), 
          multiple particles from the same event triggering different stations (in dotted yellow), 
          and outgoing neutrinos (in green).
      }
      \label{fig:coinbreakdown-allflavor}
    \end{figure}

    In the case of both tracks and decays, the tau or muon can originate inside or outside of our array's triggering region, so long as any cascades created exist within the triggering region.
    Both the standard detector model, AraSim, and \texttt{PyREx} only simulate neutrinos originating within ARA's triggering region and estimate tau and muon decays at the site of the neutrino interaction.
    While the influence of tau and muon cascades on first generation arrays is small, making this a rational estimate, larger arrays that are currently under development should use a full simulation.
    NuLeptonSim allows us to do so for this study by modeling secondary production for both muons and taus and propagation, using both continuous and stochastic energy losses.
    Neutrino generation is performed in a larger volume, outgoing taus and muons are propagated through the ice using a fast algorithm that identifies tracks that pass through a smaller volume around the array.
    Then the more time consuming signal generation can be performed for particles in a smaller volume, closer to the array~\citep{nuleptonsim}.
    NuRadioMC is a simulation framework that has a similar setting, using PROPOSAL for secondary production and propagation in a larger volume than considered for signal generation~\citep{nuradiomc1, nuradiomc2}.

    % The effective area results of our simulation are shown in Figure~\ref{fig:veffs}.
    In Figure~\ref{fig:EAbyflavor}, we compare our signal-only initial cascade results to previous ARA signal+noise simulations of initial cascades and tau/muon approximations.
    Although ARA's realistic use of noise and lepton approximations makes their effective areas larger than our calculations, our all-flavor initial cascade curve (red curve versus black curve) is reasonable.
    Note that the ARA simulation we compare to is the effective area of 1 station multiplied by 5 (to match our 5-station simulation) which overestimates the effective area at high energies since this double counts multi-station observations \citep{ara_testbed}.
    Considering we have verified our calculated effective areas, we now focus on the types and prevalence of events that make up the simulated effective area.
    
    When we incorporate outgoing taus and muons, the effective area increases considerably.
    As shown by the yellow curve in Figure~\ref{fig:EAbyparticles}, at $10^{19}$ eV, triggers on cascades from only tau tracks, muon tracks, and decays make up 30\% of ARA's effective area (20\% for those resulting from muon neutrino initial cascades, and 10\% from tau neutrinos).
    It's also clear from this plot that outgoing neutrinos from neutral current interactions or tau decays should not contribute greatly to our observations, as expected.
    From the effective areas we calculated and the 2010 Kotera et al cosmogenic star formation rate flux (with $E_{max}=10^{21.5}$eV) \citep{flux-kotera}, we expect up to 1.56 UHE neutrinos in our 25 station-year sample (0.13 via the 2019, 10\% proton, van Vliet et. al. flux \citep{flux_vanvliet}).
    The strong performance of muon and tau cascades at high energies leads to a prediction of 0.52 (0.18) charged lepton events in our 25 station-years of data, or 1 event in 48 (139) station-years of data.
    This means cascades from charged leptons make up 25\% of all triggered events, which supports findings from the 2020 simulation study published by Garc{\'\i}a-Fern{\'a}ndez et al \citep{secondaries_garciafernandez}.
    
    Figure~\ref{fig:EAbyNTriggers} shows that triggered events involving more than one triggering cascade (such as multiple energy depositions from tracks or an initial cascade and a subsequent decay) make up 13\% of triggered ARA events associated with $10^{19}$ eV initial cascades.
    An interesting event where three cascades from the same event triggered 4 out of 5 ARA stations is shown in Figure~\ref{fig:eventdisplay}.
    This event involves a triggering charged current interaction from a $10^{19}$ eV primary tau neutrino followed by two triggering cascades from the outgoing tau's track of stochastic energy depositions, one of which triggered two stations at once.
    The predicted event observation count for events where two or more particles trigger is low: 0.17 events in 25 station-years of data according to the 2010 Kotera et al model (or 1 event in 147 station-years).
    However, for events originating from $10^{19}$ eV neutrino interactions, just 62\% involve only an initial cascade from a primary neutrino interaction.
    The other 38\% of triggered events involve cascades from outgoing taus and muons showing the importance of their consideration in UHE neutrino observations, especially for larger in-ice radio arrays.
    This may be an overestimate as our simulation does not account for DAQ downtime between events, meaning signal from a cascade may reach an ARA station when that station is busy recording a previous cascade and therefore not actively taking new data.
    For triggers associated with $3\times10^{18}$ eV initial cascades the percent of triggered events including only an initial cascade is 73\% and is 85\% for $10^{18}$ eV.

  \subsection{Multi-Station Events}

    %Events that trigger more than 1 ARA station can lead to increased pointing accuracy of the Askaryan Radio Array.
    Previous studies of events that trigger more than 1 ARA station concluded that, at $10^{18}$ eV, multi-station initial cascade triggers account for 5\% of events \citep{ara_a23_10mo}.
    Our multi-station results for initial cascades (Figure~\ref{fig:coinbreakdown-allflavor}, blue curve) are consistent with the previous measurement but the multi-station trigger prevalence increases to almost 7\% when considering tracks and decays from outgoing taus and muons.
    For events where the primary neutrino has $10^{19}$ eV of energy, the total multi-station contribution increases to 17\%, where almost half of these multi-station events occur due to the involvement of muon and/or tau cascades.
    We predict there could be 0.25 multi-station events in 25 station-years of data according to the 2010 Kotera et al model.
    %(or 1 event in 100 station-years).

    Looking for events that trigger more than one station, while interesting, is a challenge.
    If we want to do fully integrated inter-station reconstruction, stations need to have synchronized clocks or that we know their relative timings.
    We attempted to synchronize ARA station clocks using a White Rabbit networking switch at the IceCube Lab, but three of our stations had incompatible SFPs (electrical-to-optical converters).
    
  \subsection{Conclusion} 

  We estimated the impact of outgoing muons, outgoing taus, and multi-station events in an array that historically has focused on only initial cascades and single-station triggering events.
  The trigger-level effective areas we calculated with our signal-only simulation predicts that the ARA 10-year dataset may contain one or two cosmogenic neutrinos.
  While tracks and multi-station events could take decades to observe in our five-station array, a 35-station array like \mbox{RNO-G} or a hundreds-of-stations array like IceCube-Gen2-Radio will be able to accumulate more station-years of data in less time and could be able to see muons, taus, and multi-station events every few years once they are fully installed.
  Not only do these arrays have more stations, but they also have next generation station hardware and designs that could allow for further increased likelihood for observing unique events.